In mathematics, it is common practice to assume a fixed set theory, usually with the axiom of choice. While there are advantages to working in one common language and sharing many of the basic constructions, the dual approach of adapting the ``base language'' to particular mathematical domains is sometimes more concise, provides a new perspective, and encourages new proof techniques which would be hard to find otherwise. We use the word ``synthetic'' to indicate that the latter approach is used.

This article is written in a language called \emph{synthetic algebraic geometry} which is introduced in \cite{draft} by adding three axioms to homotopy type theory. Although homotopy theoretic reasoning will not appear in an essential way in this article, it plays an important role in the development of the subject. We refer the reader to~\cite{draft} for an introduction to synthetic algebraic geometry and to~\cite{overview} for an updated short introduction to the subject and its history.
